
\subsection{同伦等价与算子 $\mathfrak q$ III：辛微分同胚下的不变性 2}
\label{subsec:4-6-6}
\liftIndex{I180}

本小节证明命题 4.6.55。
第 4.6.5 小节所需的
$\mathcal X_{3+}(L')$ 与 $\mathcal X'(M')$
的选取，论证类似第 7.2 节。

算子 $\mathfrak f_{\ell,k}$ 的构造
结合了第 4.6.1 和第 3.8.6 小节的构造。
沿用这两小节及第 4.6.5 小节的记号。
考虑
$(\Sigma,\vec z,\vec z^+)
\in\mathcal M_{k+1,\ell}^{\mathrm{main}}$，
并按照定义 4.6.6 对它的分支排序。
回顾：附着的球面气泡树
计入其根所在的分支。

\begin{definition}{4.6.58}\label{def:4-6-58}
定义
\[
\mathcal M_{k+1,\ell}^{\mathrm{main}}
(M',L',\{J_\rho\}_\rho:\beta;\operatorname{top}(\rho))
\]
为所有满足
(4.6.7.1)、(4.6.7.2)、(4.6.7.4)、
(4.6.7.5) 及下述条件的数据
\[
\bigl((\Sigma,\vec z,\vec z^+),
      (u_\alpha),(\rho_\alpha)\bigr)
\]
的等价类之集：

\noindent\textup{(4.6.7.3$'$)}
$((\Sigma,\vec z,\vec z^+),(u_\alpha))$
在定义 2.1.24 的意义下稳定。
\end{definition}

类似第 4.6.1 小节，定义评价映射
\[
\begin{aligned}
ev
&=(ev^{\mathrm{int}},ev_+^b,ev_0^b)\\
&=(ev_1^{\mathrm{int}},\ldots,ev_\ell^{\mathrm{int}},
   ev_1^b,\ldots,ev_k^b,ev_0^b)\\
&\colon
\mathcal M_{k+1,\ell}^{\mathrm{main}}
(M',L',\{J_\rho\}_\rho:\beta;\operatorname{top}(\rho))
\longrightarrow M^\ell\times L^k\times L'.
\end{aligned}
\]
这里上标 $b$ 表示边界，$\mathrm{int}$ 表示内部；
输入处的评价映射按第 4.6.1 小节的约定
经 $\psi^{-1}$ 取值于 $M$ 或 $L$，
输出处的 $ev_0^b$ 则取值于 $L'$。
对 $P_i\in\mathcal X_1(L)$、
$Q_i\in\mathcal X(M)$，令
\[
\begin{aligned}
&\mathcal M_{k+1,\ell}^{\mathrm{main}}
(M',L',\{J_\rho\}_\rho:\beta;\operatorname{top}(\rho);
 \vec Q,\vec P)\\
&\quad=
\mathcal M_{k+1,\ell}^{\mathrm{main}}
(M',L',\{J_\rho\}_\rho:\beta;\operatorname{top}(\rho))\\[-1mm]
&\qquad{}_{(ev^{\mathrm{int}},ev_+^b)}\!\times
\bigl(Q_1\times\cdots\times Q_\ell
      \times P_1\times\cdots\times P_k\bigr).
\end{aligned}
\]
选取适当的仓西结构与多值截面
$\mathfrak s$，并置
\[
\begin{aligned}
&\mathcal M_{k+1,\ell}^{\mathrm{main}}
(M',L',\{J_\rho\}_\rho:\beta;\operatorname{top}(\rho);
 \vec Q,\vec P)^{\mathfrak s}\\
&\qquad=
\bigl(\mathfrak s^{-1}(0),ev_0^b\bigr).
\end{aligned}
\]
右边先经三角剖分并选定顶点顺序，
视为奇异链。

\begin{definition}{4.6.59}\label{def:4-6-59}
定义
\[
\begin{aligned}
&\mathfrak f_{\ell,k;\beta}
\bigl((Q_1\otimes\cdots\otimes Q_\ell)
\otimes(P_1\otimes\cdots\otimes P_k)\bigr)\\
&\quad=
\frac{1}{\ell!}\,
\mathcal M_{k+1,\ell}^{\mathrm{main}}
(M',L',\{J_\rho\}_\rho:\beta;\operatorname{top}(\rho);
 \vec Q,\vec P)^{\mathfrak s},
\end{aligned}
\]
再置
\[
\mathfrak f_{\ell,k}
=\sum_\beta T^{\omega(\beta)}
 e^{\mu(\beta)/2}\mathfrak f_{\ell,k;\beta}.
\]
\end{definition}

于是 (4.6.57) 的证明，
直接类似定理 4.6.1 与定理 3.8.32 的证明，
留给读者。
\footnote{译者注：原书实体页 274
把此处引用写作“命题 4.6.1”；
第 4.6.1 号陈述实际为定理。
见疑误 A-0255。}
命题 4.6.55 至此证明完毕。
